\documentclass[10pt,a4paper]{amsart}
\usepackage{amsmath,amssymb,amsthm,graphicx}
\usepackage{verbatim}
\newenvironment{reference}{\comment}{\endcomment}
%\newenvironment{reference}{}{}
\newenvironment{slogan}{\comment}{\endcomment}
\newenvironment{history}{\comment}{\endcomment}

% For commutative diagrams we use Xy-pic
\usepackage[all]{xy}

% We use 2cell for 2-commutative diagrams.
\xyoption{2cell}
% Theorem environments.
%
\theoremstyle{plain}
\newtheorem{theorem}[subsection]{Theorem}
\newtheorem{proposition}[subsection]{Proposition}
\newtheorem{lemma}[subsection]{Lemma}

\theoremstyle{definition}
\newtheorem{definition}[subsection]{Definition}
\newtheorem{example}[subsection]{Example}
\newtheorem{exercise}[subsection]{Exercise}
\newtheorem{situation}[subsection]{Situation}

\theoremstyle{remark}
\newtheorem{remark}[subsection]{Remark}
\newtheorem{remarks}[subsection]{Remarks}

\numberwithin{equation}{subsection}

% Macros
%
\def\lim{\mathop{\mathrm{lim}}\nolimits}
\def\colim{\mathop{\mathrm{colim}}\nolimits}
\def\Spec{\mathop{\mathrm{Spec}}}
\def\Hom{\mathop{\mathrm{Hom}}\nolimits}
\def\Ext{\mathop{\mathrm{Ext}}\nolimits}
\def\SheafHom{\mathop{\mathcal{H}\!\mathit{om}}\nolimits}
\def\SheafExt{\mathop{\mathcal{E}\!\mathit{xt}}\nolimits}
\def\Sch{\mathit{Sch}}
\def\Mor{\mathop{\mathrm{Mor}}\nolimits}
\def\Ob{\mathop{\mathrm{Ob}}\nolimits}
\def\Sh{\mathop{\mathit{Sh}}\nolimits}
\def\NL{\mathop{N\!L}\nolimits}
\def\CH{\mathop{\mathrm{CH}}\nolimits}
\def\proetale{{pro\text{-}\acute{e}tale}}
\def\etale{{\acute{e}tale}}
\def\QCoh{\mathit{QCoh}}
\def\Ker{\mathop{\mathrm{Ker}}}
\def\Im{\mathop{\mathrm{Im}}}
\def\Coker{\mathop{\mathrm{Coker}}}
\def\Coim{\mathop{\mathrm{Coim}}}

% Boxtimes
%
\DeclareMathSymbol{\boxtimes}{\mathbin}{AMSa}{"02}

%
% Macros for moduli stacks/spaces
%
\def\QCohstack{\mathcal{QC}\!\mathit{oh}}
\def\Cohstack{\mathcal{C}\!\mathit{oh}}
\def\Spacesstack{\mathcal{S}\!\mathit{paces}}
\def\Quotfunctor{\mathrm{Quot}}
\def\Hilbfunctor{\mathrm{Hilb}}
\def\Curvesstack{\mathcal{C}\!\mathit{urves}}
\def\Polarizedstack{\mathcal{P}\!\mathit{olarized}}
\def\Complexesstack{\mathcal{C}\!\mathit{omplexes}}
% \Pic is the operator that assigns to X its picard group, usage \Pic(X)
% \Picardstack_{X/B} denotes the Picard stack of X over B
% \Picardfunctor_{X/B} denotes the Picard functor of X over B
\def\Pic{\mathop{\mathrm{Pic}}\nolimits}
\def\Picardstack{\mathcal{P}\!\mathit{ic}}
\def\Picardfunctor{\mathrm{Pic}}
\def\Deformationcategory{\mathcal{D}\!\mathit{ef}}
\title{Adjunctions between the functors on presheaves attached to an adjoint pair of functors}
\author{The Stacks Project Authors}
\date{stacks-project:sites.tex:local-snapshot-20261012; GNU FDL 1.2 or later}
\begin{document}
\maketitle
\noindent\textit{Editorial note.} All delivered mathematical text is the author's own, copied from the pinned Stacks Project chapter file sites-cohomology.tex; no statement or proof step was rewritten, reordered or removed.\par\medskip
\noindent\textit{Editorial note.} The proof is the author's complete argument; every step the author writes out is inside the counted span and no closing sentence is a placeholder.\par\medskip
\noindent\textit{Editorial note.} Every labelled reference occurring inside the counted spans is carried in assembly.external\_references as a non-load-bearing navigation pointer (same-chapter ones at their definition site, cross-chapter ones at the invocation inside the counted span), so no load-bearing mathematics is replaced by a stub. No bibliography citation occurs inside the counted spans.\par\medskip
\noindent\textit{Editorial note.} The checked chapter list for same-name collisions is the set of stacks chapters staged in this repository: chow, cohomology, divisors, etale-cohomology, flat, more-algebra, more-morphisms, morphisms, spaces-more-morphisms, varieties together with this card's chapter sites-cohomology.\par\medskip
\noindent\textit{Editorial note.} The preamble is the chapter's own retained minimal preamble (sources/stacks-sites-cohomology/retained-preamble.tex), a verbatim excerpt of the native preamble: the Xy-pic setup and the native theorem declarations with the mathematical macros.\par\medskip
\noindent\textit{Allowed standard prerequisites (editorial).}\begin{itemize}
\item The functors on presheaves attached to a functor of categories. Same chapter, the definitions of the direct and inverse image functors on presheaves and of the auxiliary categories of morphisms and of objects mapping into a fixed object. Supplies the functors and the categories appearing in the statement.
\end{itemize}\medskip
\begin{lemma}
\label{lemma-adjoint-functors}
Let $u : \mathcal{C} \to \mathcal{D}$ and $v : \mathcal{D} \to \mathcal{C}$
be functors of categories. Assume that $v$ is right adjoint to $u$.
Then we have
\begin{enumerate}
\item $u^ph_V = h_{v(V)}$ for any $V$ in $\mathcal{D}$,
\item the category $\mathcal{I}^v_U$ has an initial object,
\item the category ${}_V^u\mathcal{I}$ has a final object,
\item ${}_pu = v^p$, and
\item $u^p = v_p$.
\end{enumerate}
\end{lemma}

\begin{proof}
Proof of (1). Let $V$ be an object of $\mathcal{D}$. We have
$u^ph_V = h_{v(V)}$ because
$u^ph_V(U) = \Mor_\mathcal{D}(u(U), V) = \Mor_\mathcal{C}(U, v(V))$
by assumption.

\medskip\noindent
Proof of (2). Let $U$ be an object of $\mathcal{C}$. Let
$\eta : U \to v(u(U))$ be the map adjoint to the map
$\text{id} : u(U) \to u(U)$. Then we claim $(u(U), \eta)$
is an initial object of $\mathcal{I}_U^v$. Namely, given
an object $(V, \phi : U \to v(V))$ of $\mathcal{I}_U^v$
the morphism $\phi$ is adjoint to a map $\psi : u(U) \to V$
which then defines a morphism $(u(U), \eta) \to (V, \phi)$.

\medskip\noindent
Proof of (3). Let $V$ be an object of $\mathcal{D}$. Let
$\xi : u(v(V)) \to V$ be the map adjoint to the map
$\text{id} : v(V) \to v(V)$. Then we claim $(v(V), \xi)$
is a final object of ${}_V^u\mathcal{I}$. Namely, given
an object $(U, \psi : u(U) \to V)$ of ${}_V^u\mathcal{I}$
the morphism $\psi$ is adjoint to a map $\phi : U \to v(V)$
which then defines a morphism $(U, \psi) \to (v(V), \xi)$.

\medskip\noindent
Hence for any presheaf $\mathcal{F}$ on $\mathcal{C}$ we have
\begin{eqnarray*}
v^p\mathcal{F}(V)
& = &
\mathcal{F}(v(V)) \\
& = &
\Mor_{\textit{PSh}(\mathcal{C})}(h_{v(V)}, \mathcal{F}) \\
& = &
\Mor_{\textit{PSh}(\mathcal{C})}(u^ph_V, \mathcal{F}) \\
& = &
\Mor_{\textit{PSh}(\mathcal{D})}(h_V, {}_pu\mathcal{F}) \\
& = &
{}_pu\mathcal{F}(V)
\end{eqnarray*}
which proves part (4). Part (5) follows by the uniqueness of adjoint functors.
\end{proof}

\end{document}
